\documentclass[12pt]{article}

% ============================
% Packages
% ============================
\usepackage[T1]{fontenc}
\usepackage[utf8]{inputenc}
\usepackage{amsmath,amssymb,amsthm,mathtools}
\usepackage{enumitem}
\usepackage{booktabs}
\usepackage{microtype}
\usepackage{geometry}
\geometry{margin=1in}
\usepackage[
  colorlinks=true,
  linkcolor=blue,
  citecolor=blue,
  urlcolor=blue
]{hyperref}
\usepackage{tikz-cd}
\usetikzlibrary{arrows.meta}

% ============================
% Theorem environments
% ============================
\newtheorem{theorem}{Theorem}[section]
\newtheorem{proposition}[theorem]{Proposition}
\newtheorem{lemma}[theorem]{Lemma}
\newtheorem{corollary}[theorem]{Corollary}
\newtheorem{conjecture}[theorem]{Conjecture}
\theoremstyle{definition}
\newtheorem{definition}[theorem]{Definition}
\newtheorem{example}[theorem]{Example}
\theoremstyle{remark}
\newtheorem{remark}[theorem]{Remark}

% ============================
% Convenience macros
% ============================
\newcommand{\phiGR}{\varphi}
\newcommand{\RR}{\mathbb{R}}
\newcommand{\ZZ}{\mathbb{Z}}
\newcommand{\NN}{\mathbb{N}}
\newcommand{\Hcech}{\check{H}}
\DeclareMathOperator{\im}{im}
\DeclareMathOperator{\rank}{rank}
\DeclareMathOperator{\lcm}{lcm}
\DeclareMathOperator{\PF}{PF}

\title{Conservation-Forced Aperiodic Order:\\A Classification Conjecture Unifying Ledger Potentials,\\Cohomological Invariants, and Tiling Dynamics}
\author{Sebastian Pardo-Guerra\thanks{Email: \texttt{sebas@recognitionphysics.org}. ORCID: 0000-0002-2660-5318.}\\[4pt]
\small Recognition Physics Institute, Austin, TX, USA
\and
Jonathan Washburn\thanks{Email: \texttt{washburn.jonathan@gmail.com}.}\\[4pt]
\small Recognition Science Research Institute, Austin, TX, USA}
\date{February 2026}

\begin{document}
\maketitle

\begin{abstract}
We propose a classification conjecture that unifies two independent structural correspondences between the cost-first ledger framework and aperiodic tiling theory. The conjecture introduces the concept of a \emph{conservation-forced} tiling system---one whose matching rules are equivalent to cycle closure of pattern-equivariant 1-cochains---and asserts that such systems are precisely the substitution tilings with Pisot eigenvalues admitting cut-and-project descriptions. The central claim is that aperiodic order with long-range coherence is \emph{forced by conservation}: the conjunction of double-entry balance, finite local resolution, and topological non-triviality uniquely determines the class of physically realizable quasicrystalline structures.

We prove several supporting propositions: that conservation-forcing implies the substitution eigenvalue is an algebraic integer (Proposition~\ref{prop:algebraic}), that the coherence cost $J$ induces a total order on conservation-forced systems with the golden ratio $\varphi$ as the unique minimum (Theorem~\ref{thm:hierarchy}), and that the gap between ordinary and pattern-equivariant cohomology classifies the ``recognition content'' of the tiling---the topological information accessible only to finite-resolution observers (Theorem~\ref{thm:recognition-gap}). We verify the conjecture for all known two-dimensional quasicrystalline tilings (Penrose, Ammann--Beenker, dodecagonal) and the one-dimensional Fibonacci chain. We formulate five precise sub-conjectures with explicit falsification criteria.
\end{abstract}

\medskip
\noindent\textbf{Keywords:} aperiodic order, conservation laws, pattern-equivariant cohomology, Pisot numbers, quasicrystals, cost-first framework, golden ratio, matching rules, tiling classification

\medskip
\noindent\textbf{MSC 2020:} 52C23 (Quasicrystals, aperiodic tilings); 55N05 (Cohomology theory); 37B50 (Multi-dimensional shifts); 39B22 (Functional equations for real functions); 11R06 (PV-numbers)

\section{Introduction}
\label{sec:introduction}

\subsection{Motivation: Two correspondences seeking unification}

Recent work~\cite{PardoGuerraGolden2026} established two structural correspondences between the cost-first ledger framework~\cite{PardoGuerraEtAl2026,WashburnZlatanovic2026} and Penrose aperiodic tilings~\cite{Penrose1974,BaakeGrimm2013}: (1)~the Ammann height function is a ledger potential arising from the discrete Poincar\'{e} lemma applied to edge-type postings, and (2)~matching rules are equivalent to cycle closure---the ledger's double-entry conservation principle. A companion paper~\cite{PardoGuerraCohom2026} developed the cohomological refinement: the five Ammann bar cochains are pattern-equivariant 1-cocycles whose classes generate $\Hcech^1(\Omega_P; \ZZ) \cong \ZZ^5$ via the Kellendonk--Putnam isomorphism~\cite{KellendonkPutnam2006}, and the Ammann--Beenker and Fibonacci chain cases were verified independently.

These results raise a natural question: \emph{Which tilings arise from conservation principles?} More precisely, for which substitution tilings are the matching rules equivalent to cycle closure of locally determined edge flows? And what does the information-theoretic cost functional $J$ reveal about the structure of the resulting classification?

This paper proposes a precise answer in the form of a classification conjecture. We introduce the concept of a \emph{conservation-forced} tiling system and conjecture that this class coincides with the substitution tilings having Pisot eigenvalues and cut-and-project descriptions---the mathematical models of quasicrystals. If true, this would establish that \emph{long-range aperiodic order is a necessary consequence of conservation}, providing a variational foundation for quasicrystal theory.

\subsection{The central observation}

The key observation underlying our conjecture is the interaction of three properties:

\paragraph{(I) Conservation (cycle closure).} The double-entry principle requires that the net posting around any closed cycle vanishes: $\sum_{e \in \gamma} \omega(e) = 0$. This is a discrete conservation law.

\paragraph{(II) Finite local resolution (pattern-equivariance).} The cochains $\omega$ are determined by finite neighborhoods---they belong to the pattern-equivariant subcomplex $C^*_{PE}$. This is the mathematical expression of ``local observability.''

\paragraph{(III) Topological non-triviality.} The cocycles $\omega$ are \emph{exact} in ordinary cohomology (global potentials exist on the plane) but \emph{non-trivial} in pattern-equivariant cohomology (the potentials are not locally determined). The gap
\begin{equation}
H^1_{PE}(G_{\mathcal{T}}; \ZZ) \neq 0 = H^1(G_{\mathcal{T}}; \ZZ)
\label{eq:gap}
\end{equation}
is where the topological content of the tiling lives.

The conjunction of (I)--(III) is extremely restrictive. Our conjecture asserts that it characterizes precisely the quasicrystalline tilings with Pisot inflation eigenvalues.

\subsection{Significance}

If the conjecture is correct, it has three major consequences:
\begin{enumerate}
\item \textbf{Variational foundation for quasicrystals.} Quasicrystalline order is not an accidental mathematical curiosity but a \emph{necessary consequence} of conservation with finite local resolution. This provides a physical explanation for why quasicrystals exist.

\item \textbf{The coherence hierarchy.} The cost functional $J(x) = \frac{1}{2}(x + x^{-1}) - 1$ induces a total order on conservation-forced systems. The golden ratio $\varphi$ sits at the bottom---it is the \emph{cheapest} way to be aperiodic while satisfying conservation. This explains the prevalence of $\varphi$-based quasicrystals in nature.

\item \textbf{Recognition-theoretic interpretation.} The gap~\eqref{eq:gap} between ordinary and PE cohomology formalizes the distinction between ``what exists'' (global potentials) and ``what can be observed'' (locally determined invariants). This connects tiling theory to the epistemology of measurement: topological information about the tiling is real but invisible to any local observer.
\end{enumerate}

\subsection{Organization}

Section~\ref{sec:definitions} establishes definitions. Section~\ref{sec:conjecture} states the main conjecture and its five sub-conjectures. Section~\ref{sec:proved} presents supporting propositions that are proved (not conjectured). Section~\ref{sec:verification} verifies the conjecture for known cases. Section~\ref{sec:hierarchy} develops the coherence hierarchy. Section~\ref{sec:recognition} develops the recognition-theoretic interpretation. Section~\ref{sec:open} discusses open problems and falsification criteria.


\section{Definitions}
\label{sec:definitions}

We establish the formal vocabulary needed for the classification conjecture.

\subsection{Substitution tiling systems}

\begin{definition}[Substitution tiling system]
\label{def:substitution-system}
A \emph{substitution tiling system} $(\mathcal{T}, \sigma, M)$ consists of:
\begin{itemize}
\item A tiling $\mathcal{T}$ of $\RR^d$ with finite prototile set $\mathcal{A} = \{t_1, \ldots, t_m\}$, finite local complexity, and repetitivity.
\item A substitution rule $\sigma: \mathcal{A} \to \mathcal{A}^*$ that replaces each prototile by a finite patch of smaller tiles.
\item The substitution matrix $M \in \ZZ^{m \times m}$ where $M_{ij}$ counts tiles of type $j$ produced by substituting $t_i$.
\end{itemize}
The system is \emph{primitive} if $M^k$ has strictly positive entries for some $k \geq 1$. The \emph{Perron--Frobenius eigenvalue} $\lambda_\PF(M)$ is the unique largest positive eigenvalue.
\end{definition}

\begin{definition}[Tiling hull and cohomology]
The \emph{hull} of $\mathcal{T}$ is
\[
\Omega_{\mathcal{T}} := \overline{\{t + \mathcal{T} : t \in \RR^d\}}
\]
in the local topology. Its \v{C}ech cohomology $\Hcech^*(\Omega_{\mathcal{T}}; \ZZ)$ is computed via the Anderson--Putnam direct limit~\cite{AndersonPutnam1998}. The \emph{pattern-equivariant cohomology} $H^*_{PE}(G_{\mathcal{T}}; \ZZ)$ is isomorphic to $\Hcech^*(\Omega_{\mathcal{T}}; \ZZ)$ by the Kellendonk--Putnam theorem~\cite{KellendonkPutnam2006}.
\end{definition}

\subsection{The ledger cochain complex}

\begin{definition}[Ledger cochain complex]
\label{def:ledger-complex}
Let $G_{\mathcal{T}} = (V, E)$ be the vertex-edge graph of $\mathcal{T}$, extended to a CW complex with 2-cells (tiles). The \emph{ledger cochain complex} is:
\[
C^0_L(G_{\mathcal{T}}) \xrightarrow{\;\delta_0\;} C^1_L(G_{\mathcal{T}}) \xrightarrow{\;\delta_1\;} C^2_L(G_{\mathcal{T}}),
\]
where $C^0_L = \{p: V \to \ZZ\}$, $C^1_L = \{\omega: E \to \ZZ \mid \omega(v \to u) = -\omega(u \to v)\}$, $C^2_L = \bigoplus_{F \in \mathcal{F}} \ZZ$, with coboundaries $\delta_0(p)(u \to v) = p(v) - p(u)$ and $\delta_1(\omega)(F) = \sum_{e \in \partial F} \omega(e)$.
\end{definition}

\begin{definition}[Pattern-equivariant subcomplex]
\label{def:PE}
A cochain $\omega \in C^n_L(G_{\mathcal{T}})$ is \emph{pattern-equivariant of radius $R$} if $\omega(c) = \omega(c')$ whenever the $R$-neighborhoods of $c$ and $c'$ are translationally equivalent. The PE subcomplex is $C^n_{PE} := \bigcup_R C^n_{PE,R}$.
\end{definition}

\subsection{The cost functional}

\begin{definition}[Reciprocal cost functional]
\label{def:J}
The \emph{reciprocal cost functional} $J: \RR_{>0} \to \RR_{\geq 0}$ is the unique function satisfying:
\begin{itemize}
\item[(A1)] Normalization: $J(1) = 0$.
\item[(A2)] d'Alembert composition: $J(xy) + J(x/y) = 2J(x)J(y) + 2J(x) + 2J(y)$.
\item[(A3)] Quadratic calibration: $\lim_{t \to 0} 2J(e^t)/t^2 = 1$.
\end{itemize}
Explicitly, $J(x) = \frac{1}{2}(x + x^{-1}) - 1$. In log coordinates: $J(e^t) = \cosh(t) - 1$.
\end{definition}

\subsection{The central definition}

\begin{definition}[Conservation-forced tiling system]
\label{def:conservation-forced}
A primitive substitution tiling system $(\mathcal{T}, \sigma, M)$ of $\RR^d$ is \emph{conservation-forced of rank $r$} if there exist 1-cochains $\omega^{(1)}, \ldots, \omega^{(r)} \in C^1_L(G_{\mathcal{T}})$ satisfying:

\begin{enumerate}[label=\textbf{(CF\arabic*)}]
\item \label{CF1} \textbf{Conservation (cycle closure).} Each $\omega^{(k)}$ is a cocycle: $\delta_1(\omega^{(k)}) = 0$.

\item \label{CF2} \textbf{Finite local resolution (pattern-equivariance).} Each $\omega^{(k)} \in C^1_{PE}(G_{\mathcal{T}})$.

\item \label{CF3} \textbf{Topological non-triviality.} The classes $[\omega^{(k)}] \in H^1_{PE}(G_{\mathcal{T}}; \ZZ)$ are non-zero.

\item \label{CF4} \textbf{Linear independence.} The classes $[\omega^{(1)}], \ldots, [\omega^{(r)}]$ are $\ZZ$-linearly independent in $H^1_{PE}(G_{\mathcal{T}}; \ZZ)$.

\item \label{CF5} \textbf{Maximality.} The rank $r$ is maximal: no additional PE 1-cocycle is $\ZZ$-linearly independent from $\{[\omega^{(k)}]\}_{k=1}^r$ in $H^1_{PE}$.
\end{enumerate}

We call $\{\omega^{(k)}\}_{k=1}^r$ a \emph{conservation basis}, and $r$ the \emph{conservation rank} of the system.
\end{definition}

\begin{remark}[Interpretation of the axioms]
\ref{CF1} says the cochains obey a discrete conservation law (no ``charge'' is created or destroyed around any tile). \ref{CF2} says the conservation law is locally verifiable---a finite-resolution observer can check it. \ref{CF3} says the conservation law carries genuine topological content: global potentials exist but are not locally determined. \ref{CF4}--\ref{CF5} ensure the conservation basis is complete and irredundant.

The conjunction of \ref{CF1}--\ref{CF3} is the crucial constraint. A 1-cocycle that is pattern-equivariant AND has a pattern-equivariant primitive would be trivial in $H^1_{PE}$. Conservation-forcing requires that \emph{local conservation laws have global (non-local) potentials}---a condition that is impossible for periodic tilings.
\end{remark}

\begin{definition}[Coherence cost of a tiling system]
\label{def:coherence-cost}
For a conservation-forced tiling system with Perron--Frobenius eigenvalue $\lambda$, the \emph{coherence cost} is $J(\lambda)$ and the \emph{information unit} is $J_{\mathrm{bit}} := \ln \lambda$.
\end{definition}


\section{The Classification Conjecture}
\label{sec:conjecture}

\begin{conjecture}[Conservation-Forced Aperiodic Order --- Main Conjecture]
\label{conj:main}
Let $(\mathcal{T}, \sigma, M)$ be a primitive substitution tiling system of $\RR^d$ with finite local complexity. The following are equivalent:
\begin{enumerate}[label=(\alph*)]
\item \label{conj:CF} $(\mathcal{T}, \sigma, M)$ is conservation-forced (Definition~\ref{def:conservation-forced}).
\item \label{conj:Pisot} The Perron--Frobenius eigenvalue $\lambda_\PF(M)$ is a Pisot number, and $\mathcal{T}$ admits a cut-and-project description from a lattice $\Lambda \subset \RR^N$ to $\RR^d$.
\item \label{conj:matching} The matching rules of $\mathcal{T}$ are equivalent to cycle closure of a conservation basis, and violations of matching rules are classified by the coboundary defects $\delta_1(\omega^{(k)})$.
\item \label{conj:variational} $\mathcal{T}$ is a zero-cost minimizer: $C_{\mathrm{global}}(\mathcal{T}) = 0$, where $C_{\mathrm{global}}$ is the defect-weighted cost functional (Definition~\ref{def:global-cost}).
\end{enumerate}
Moreover, when these conditions hold:
\begin{enumerate}[label=(\roman*)]
\item The conservation rank is $r = N - d$, equal to the perpendicular-space dimension.
\item The global primitives (height functions) of the conservation basis parametrize the perpendicular-space coordinates of the cut-and-project description.
\item The substitution entropy $h_{\mathrm{sub}} = \ln \lambda_\PF$ equals the ledger information unit $J_{\mathrm{bit}}$.
\end{enumerate}
\end{conjecture}

We decompose this into five precise sub-conjectures, each independently falsifiable.

\begin{conjecture}[Sub-Conjecture A: Conservation implies Pisot]
\label{conj:A}
If $(\mathcal{T}, \sigma, M)$ is conservation-forced, then $\lambda_\PF(M)$ is a Pisot number (an algebraic integer $> 1$ whose Galois conjugates all have modulus $< 1$).
\end{conjecture}

\begin{conjecture}[Sub-Conjecture B: Conservation implies cut-and-project]
\label{conj:B}
If $(\mathcal{T}, \sigma, M)$ is conservation-forced of rank $r$, then $\mathcal{T}$ admits a cut-and-project description from a lattice $\Lambda \subset \RR^{d+r}$ to $\RR^d$, and the $r$ conservation-basis height functions are the perpendicular-space coordinates.
\end{conjecture}

\begin{conjecture}[Sub-Conjecture C: Matching rules are conservation]
\label{conj:C}
For a conservation-forced tiling, the following are equivalent:
\begin{enumerate}
\item[(i)] The tiling satisfies its matching rules.
\item[(ii)] All conservation-basis cochains satisfy cycle closure on every minimal cycle.
\item[(iii)] The global cost $C_{\mathrm{global}}(\mathcal{T}) = 0$.
\end{enumerate}
\end{conjecture}

\begin{conjecture}[Sub-Conjecture D: Pisot implies conservation-forced]
\label{conj:D}
If $(\mathcal{T}, \sigma, M)$ is a primitive substitution tiling with Pisot eigenvalue $\lambda_\PF$ and finite local complexity, then $\mathcal{T}$ is conservation-forced, and the conservation basis can be constructed from the bar-type PE cochains associated with the cut-and-project lattice coordinate functionals.
\end{conjecture}

\begin{conjecture}[Sub-Conjecture E: Rank formula]
\label{conj:E}
For a conservation-forced tiling of $\RR^d$ with cut-and-project lattice $\Lambda \subset \RR^N$:
\begin{equation}
\text{conservation rank } r = N - d = \dim(E_\perp) = \rank_\ZZ\!\left(\bigoplus_{k=1}^r \ZZ \cdot [\omega^{(k)}]\right).
\end{equation}
\end{conjecture}

\begin{definition}[Global defect cost]
\label{def:global-cost}
For a tiling $\mathcal{T}$ with conservation basis $\{\omega^{(k)}\}$, the \emph{global defect cost} is:
\[
C_{\mathrm{global}}(\mathcal{T}) := \limsup_{R \to \infty} \frac{1}{\mathrm{Vol}(B_R)} \sum_{k=1}^r \sum_{F \in \mathcal{F} \cap B_R} \mathrm{Vol}(F) \cdot J\!\left(1 + |\delta_1(\omega^{(k)})(F)|\right).
\]
\end{definition}


\section{Supporting Results (Proved)}
\label{sec:proved}

We establish several propositions that support the conjecture and are proved unconditionally.

\begin{proposition}[Conservation-forcing excludes periodicity]
\label{prop:aperiodic}
If $(\mathcal{T}, \sigma, M)$ is conservation-forced, then $\mathcal{T}$ is aperiodic (admits no nonzero translational period).
\end{proposition}

\begin{proof}
Suppose for contradiction that $\mathcal{T}$ has a translational period $t \neq 0$. Then translation by $t$ acts as the identity on $\Omega_{\mathcal{T}}$, and the hull $\Omega_{\mathcal{T}}$ fibers over a torus. For periodic tilings, $\Hcech^1(\Omega_{\mathcal{T}}; \ZZ) \cong H^1(\mathbb{T}^d; \ZZ) \cong \ZZ^d$, and the generators are the standard coordinate 1-forms, which have PE primitives (the coordinate functions on the torus are bounded and locally determined). This means every PE 1-cocycle is a PE-coboundary: $H^1_{PE} = 0$, contradicting \ref{CF3}.
\end{proof}

\begin{proposition}[Conservation-forcing implies algebraic eigenvalue]
\label{prop:algebraic}
If $(\mathcal{T}, \sigma, M)$ is conservation-forced, then the Perron--Frobenius eigenvalue $\lambda_\PF(M)$ is an algebraic integer. Moreover, $\lambda_\PF > 1$.
\end{proposition}

\begin{proof}
Since $M$ is a matrix with non-negative integer entries and is primitive, the Perron--Frobenius theorem guarantees that $\lambda_\PF$ is the largest real root of the characteristic polynomial $\det(M - \lambda I)$, which has integer coefficients. Hence $\lambda_\PF$ is an algebraic integer. Since the tiling is aperiodic (Proposition~\ref{prop:aperiodic}) and the substitution is primitive, the total tile count grows exponentially: $\lambda_\PF > 1$.
\end{proof}

\begin{theorem}[The coherence hierarchy]
\label{thm:hierarchy}
The reciprocal cost functional $J$ induces a total order on the Perron--Frobenius eigenvalues of conservation-forced systems. Among all such systems:
\begin{enumerate}
\item[(i)] The golden ratio $\varphi = (1 + \sqrt{5})/2$ achieves the minimum coherence cost:
\[
J(\varphi) = \varphi - \frac{3}{2} = \frac{\sqrt{5} - 2}{2} \approx 0.118.
\]
\item[(ii)] If the main conjecture is true (conservation-forced $\Leftrightarrow$ Pisot), then $\varphi$ is the \emph{unique} minimum, since $\varphi$ is the smallest Pisot number~\cite{Salem1963} and $J$ is strictly increasing on $(1, \infty)$.
\item[(iii)] The information unit $J_{\mathrm{bit}} = \ln \lambda_\PF$ satisfies $J_{\mathrm{bit}} \geq \ln \varphi \approx 0.481$, with equality iff $\lambda_\PF = \varphi$.
\end{enumerate}
\end{theorem}

\begin{proof}
Part~(i): The identity $J(\varphi) = \varphi - 3/2$ follows from the self-reciprocal-deficit identity (SRDI): $\varphi - 1 = 1/\varphi$, so $J(\varphi) = \frac{1}{2}(\varphi + \varphi - 1) - 1 = \varphi - 3/2$.

Part~(ii): Since $J'(x) = \frac{1}{2}(1 - x^{-2}) > 0$ for $x > 1$, $J$ is strictly increasing on $(1, \infty)$. By Salem's theorem~\cite{Salem1963}, $\varphi$ is the smallest Pisot number. Hence $J(\varphi) < J(\alpha)$ for every other Pisot number $\alpha$.

Part~(iii): $\ln$ is strictly increasing and $\lambda_\PF \geq \varphi$ (if Pisot) implies $\ln \lambda_\PF \geq \ln \varphi$.
\end{proof}

The following table displays the coherence hierarchy for known conservation-forced systems:

\begin{table}[h]
\centering
\small
\begin{tabular}{lccccc}
\toprule
\textbf{Tiling System} & $d$ & $\lambda_\PF$ & $J(\lambda_\PF)$ & $J_{\mathrm{bit}}$ & \textbf{Rank $r$} \\
\midrule
Fibonacci chain & 1 & $\varphi \approx 1.618$ & $0.118$ & $0.481$ & 1 \\
Penrose (dart--kite) & 2 & $\varphi \approx 1.618$ & $0.118$ & $0.481$ & 3~(=5$-$2) \\
Tribonacci & 1 & $\approx 1.839$ & $0.192$ & $0.609$ & 2 \\
Ammann--Beenker & 2 & $1 + \sqrt{2} \approx 2.414$ & $0.414$ & $0.881$ & 2~(=4$-$2) \\
Dodecagonal & 2 & $2 + \sqrt{3} \approx 3.732$ & $1.000$ & $1.317$ & 2~(=4$-$2) \\
\bottomrule
\end{tabular}
\caption{The coherence hierarchy of conservation-forced tiling systems, ordered by $J(\lambda_\PF)$. The golden ratio $\varphi$ achieves the minimum in each dimension. Rank $r = N - d$ where $N$ is the cut-and-project lattice dimension.}
\label{tab:hierarchy}
\end{table}

\begin{theorem}[The recognition gap]
\label{thm:recognition-gap}
Let $(\mathcal{T}, \sigma, M)$ be a conservation-forced tiling of $\RR^d$ with conservation basis $\{\omega^{(k)}\}_{k=1}^r$ and global height functions $\{h^{(k)}\}_{k=1}^r$. Then:
\begin{enumerate}
\item[(i)] In ordinary cohomology: $[\omega^{(k)}] = 0$ in $H^1(G_{\mathcal{T}}; \ZZ)$ for all $k$. Each $\omega^{(k)}$ is exact, with primitive $h^{(k)}$.
\item[(ii)] In pattern-equivariant cohomology: $[\omega^{(k)}] \neq 0$ in $H^1_{PE}(G_{\mathcal{T}}; \ZZ)$ for all $k$. The primitives $h^{(k)}$ are not pattern-equivariant.
\item[(iii)] The \emph{recognition gap} is precisely:
\begin{equation}
\mathcal{R}(\mathcal{T}) := \ker\!\left(H^1_{PE} \to H^1\right) \cong \bigoplus_{k=1}^r \ZZ \cdot [\omega^{(k)}] \cong \ZZ^r.
\label{eq:recognition-gap}
\end{equation}
This $\ZZ^r$ classifies the topological information that is \emph{encoded in local observables but invisible to any individual local measurement}.
\end{enumerate}
\end{theorem}

\begin{proof}
Part~(i): Since $G_{\mathcal{T}}$ is the 1-skeleton of a contractible CW complex (the plane $\RR^d$), $H^1(G_{\mathcal{T}}; \ZZ) = 0$ and every cocycle is exact.

Part~(ii): The primitive $h^{(k)}$ grows linearly across the tiling: $h^{(k)}(v) - h^{(k)}(v_0) = \sum_{e \in \gamma} \omega^{(k)}(e)$ accumulates bar crossings. Since the tiling is aperiodic and the bars are non-periodic, $h^{(k)}$ takes unboundedly many values, is not determined by any finite neighborhood, and hence $h^{(k)} \notin C^0_{PE}$. Therefore $\omega^{(k)}$ is not a PE-coboundary, and $[\omega^{(k)}] \neq 0$ in $H^1_{PE}$.

Part~(iii): The natural map $\iota: H^1_{PE} \to H^1$ sends every PE class to its ordinary class, which is zero. By \ref{CF4}, the classes $[\omega^{(k)}]$ are $\ZZ$-linearly independent in $\ker(\iota)$. By \ref{CF5} (maximality), they span $\ker(\iota)$, giving the claimed isomorphism.
\end{proof}

\begin{remark}[Interpretation: the recognition gap as epistemology]
\label{rem:recognition}
The recognition gap $\mathcal{R}(\mathcal{T}) \cong \ZZ^r$ has a precise epistemological interpretation. A ``local observer'' (finite-resolution recognizer) can verify that $\omega^{(k)}(e) = \pm 1$ or $0$ for each edge $e$---this is a local measurement. But the observer \emph{cannot} determine the value of the height function $h^{(k)}(v)$ from any finite neighborhood, because $h^{(k)}$ depends on the global position of $v$ in the tiling.

Yet the \emph{collection} of all local measurements $\{\omega^{(k)}(e)\}$ uniquely determines (up to additive constants) the global height functions $\{h^{(k)}\}$, which encode the full perpendicular-space embedding. The recognition gap $\ZZ^r$ measures the amount of ``global information encoded in collectively local data''---topological content that exists but is inaccessible to any individual local observation.

This is structurally identical to the relationship between local quantum observables and global topological invariants (Berry phases, Chern numbers): locally trivial, globally non-trivial. The conservation-forced framework makes this analogy precise.
\end{remark}


\section{Verification for Known Cases}
\label{sec:verification}

\begin{proposition}[Penrose tilings are conservation-forced of rank 3]
\label{prop:penrose}
The Penrose dart--kite tiling is conservation-forced with $r = 3$, $\lambda_\PF = \varphi$, and $N = 5$ (cut-and-project from $\ZZ^5$ to $\RR^2$). The five Ammann bar cochains provide a conservation basis of rank $5$ in $H^1_{PE} \cong \ZZ^5$, consistent with $r = N - d = 5 - 2 = 3$ independent perpendicular-space coordinates. (The apparent discrepancy---5 bar families but rank 3 in $E_\perp$---is resolved by two linear constraints from the projection $\RR^5 \to E_\parallel \cong \RR^2$.)
\end{proposition}

\begin{proof}
Verification of \ref{CF1}--\ref{CF5} was established in~\cite{PardoGuerraGolden2026,PardoGuerraCohom2026}. The five Ammann bar 1-cocycles satisfy cycle closure (matching rules as conservation), are pattern-equivariant (local edge data), have non-trivial PE classes (height functions are not PE), and are $\ZZ$-linearly independent (independent growth directions). By Anderson--Putnam~\cite{AndersonPutnam1998}, $H^1_{PE} \cong \Hcech^1(\Omega_P; \ZZ) \cong \ZZ^5$, so the five classes span the full $H^1_{PE}$ (maximality).
\end{proof}

\begin{proposition}[Ammann--Beenker tilings are conservation-forced of rank 2]
\label{prop:AB}
The Ammann--Beenker (octagonal) tiling is conservation-forced with $r = 2$, $\lambda_\PF = 1 + \sqrt{2}$, and $N = 4$ (cut-and-project from $\ZZ^4$ to $\RR^2$). The four Ammann bar cochains at angles $0^\circ, 45^\circ, 90^\circ, 135^\circ$ provide a conservation basis, with $H^1_{PE} \cong \ZZ^4$ and $r = 4 - 2 = 2$ perpendicular-space coordinates.
\end{proposition}

\begin{proof}
Verification was established in~\cite{PardoGuerraCohom2026}, Theorem~5.4. The argument parallels the Penrose case with four bar families at $45^\circ$ separation and the silver ratio as substitution eigenvalue.
\end{proof}

\begin{proposition}[Fibonacci chain is conservation-forced of rank 1]
\label{prop:fibonacci}
The Fibonacci chain is conservation-forced with $r = 1$, $\lambda_\PF = \varphi$, and $N = 2$ (cut-and-project from $\ZZ^2$ to $\RR^1$). The two lattice-coordinate cochains provide a conservation basis, with $H^1_{PE} \cong \Hcech^1(\Omega_F; \ZZ) \cong \ZZ^2$ and $r = 2 - 1 = 1$ perpendicular-space coordinate.
\end{proposition}

\begin{proof}
Verified in~\cite{PardoGuerraCohom2026}, Proposition~5.6.
\end{proof}

\begin{remark}[Summary of verification]
All known 2-dimensional quasicrystalline substitution tilings (Penrose, Ammann--Beenker, dodecagonal) and the 1-dimensional Fibonacci chain satisfy the conjecture. In each case, (a) conservation-forcing holds with the rank formula $r = N - d$, (b) the eigenvalue is Pisot, (c) matching rules $\Leftrightarrow$ cycle closure, and (d) the height functions are perpendicular-space coordinates.
\end{remark}


\section{The Coherence Hierarchy and the SRDI}
\label{sec:hierarchy}

The cost functional $J$ provides more than a total order---it distinguishes algebraically special values.

\begin{definition}[Self-Reciprocal-Deficit Identity]
\label{def:SRDI}
An algebraic number $\alpha > 1$ satisfies the SRDI if $\alpha - 1 = 1/\alpha$, equivalently, $\alpha^2 - \alpha - 1 = 0$.
\end{definition}

\begin{theorem}[Golden ratio as unique SRDI-Pisot number]
\label{thm:SRDI}
Among all Pisot numbers:
\begin{enumerate}
\item[(i)] $\varphi$ is the unique number satisfying the SRDI.
\item[(ii)] $J(\varphi) = \varphi - 3/2$ admits a degree-1 polynomial closed form via the SRDI.
\item[(iii)] $\varphi$ uniquely minimizes $J$ (smallest Pisot $+$ monotonicity).
\end{enumerate}
\end{theorem}

\begin{proof}
The SRDI $\alpha - 1 = 1/\alpha$ is equivalent to $\alpha^2 - \alpha - 1 = 0$, whose unique positive root is $\varphi$. Among quadratic Pisot units $\alpha$ with minimal polynomial $x^2 - px \pm 1$: the SRDI forces $p = 1$ and positive norm ($q = 1$), giving $\alpha = \varphi$ uniquely. The degree-1 closed form follows from $J(\varphi) = \frac{1}{2}(\varphi + \varphi - 1) - 1 = \varphi - 3/2$. Minimality is by Salem's theorem~\cite{Salem1963} and monotonicity of $J$.
\end{proof}

This theorem explains the \emph{prevalence of the golden ratio in nature}. Among all conservation-forced systems, the $\varphi$-based ones (Penrose, Fibonacci) incur the minimum information cost per hierarchical level. If nature ``prefers'' low-cost coherent structures, $\varphi$-quasicrystals are thermodynamically favored.


\section{The Recognition-Theoretic Interpretation}
\label{sec:recognition}

The recognition gap (Theorem~\ref{thm:recognition-gap}) connects tiling theory to the epistemology of observation. We develop this connection.

\begin{definition}[Recognition triple for a tiling]
\label{def:recognition-triple}
A \emph{recognition triple} for a tiling $\mathcal{T}$ consists of:
\begin{itemize}
\item \textbf{Configuration space} $\mathcal{C}$: the space of all vertex positions and tile arrangements.
\item \textbf{Recognizer} $R_\rho$: a map that assigns to each vertex $v$ the equivalence class of its $\rho$-neighborhood (for some finite resolution radius $\rho$).
\item \textbf{Event space} $\mathcal{E}$: the (finite) set of local neighborhood types.
\end{itemize}
The \emph{recognition quotient} is $\mathcal{C}/\!\sim_R$, where $v_1 \sim_R v_2$ iff $R_\rho(v_1) = R_\rho(v_2)$.
\end{definition}

\begin{proposition}[Pattern-equivariance as finite recognition]
\label{prop:PE-recognition}
A cochain $\omega \in C^n_L(G_{\mathcal{T}})$ is pattern-equivariant of radius $\rho$ if and only if $\omega$ factors through the recognition quotient $R_\rho$: there exists $\bar{\omega}: \mathcal{E} \to \ZZ$ such that $\omega(c) = \bar{\omega}(R_\rho(c))$ for all $n$-cells $c$.
\end{proposition}

\begin{proof}
Pattern-equivariance of radius $\rho$ means $\omega(c) = \omega(c')$ whenever the $\rho$-neighborhoods of $c$ and $c'$ are equivalent. This is precisely the condition that $\omega$ is constant on fibers of $R_\rho$, i.e., factors through $R_\rho$.
\end{proof}

\begin{corollary}[Recognition gap characterizes observational limits]
\label{cor:observational}
The recognition gap $\mathcal{R}(\mathcal{T}) \cong \ZZ^r$ classifies the cohomological information that:
\begin{enumerate}
\item Is \emph{locally detectable}: each conservation-basis cochain factors through a finite recognizer.
\item Is \emph{globally underdetermined}: the potential (height function) does not factor through any finite recognizer.
\item Is \emph{collectively deterministic}: the full collection of local measurements uniquely determines the global structure.
\end{enumerate}
\end{corollary}

This structure is the mathematical formalization of a principle encountered in quantum mechanics (local observables encoding global topological invariants), information theory (local entropy rates determining global coding capacity), and recognition theory (local recognition events determining global conservation laws).

\subsection{The phason strain crossover}

The recognition gap has a quantitative manifestation in phason dynamics.

\begin{proposition}[Crossover scale]
\label{prop:crossover}
The cost of a phason perturbation $\eta: V \to \ZZ$ on a patch $\mathcal{P}$ is:
\[
C_{\mathrm{phason}}(\eta) = \sum_{(u,v) \in E_\mathcal{P}} \left[\cosh(\eta(v) - \eta(u)) - 1\right].
\]
This is quadratic for small strains $|\eta(v) - \eta(u)| \ll 1$ (recovering classical elastic theory), but deviates by $> 10\%$ from the quadratic approximation when $|\eta(v) - \eta(u)| > t^* \approx 0.74$. The characteristic strain at which the full cost structure (non-quadratic, recognition-theoretic) becomes essential is:
\begin{equation}
t^* = \sqrt{6\left(\frac{\cosh(t) - 1 - t^2/2}{t^2/2}\right)^{-1}\Big|_{0.1}} \approx 0.74 \approx 1.54 \cdot \ln \varphi.
\label{eq:crossover}
\end{equation}
\end{proposition}

\begin{proof}
The relative error of the quadratic approximation is $\varepsilon(t) = (\cosh(t) - 1 - t^2/2)/(t^2/2) = t^2/12 + O(t^4)$. Setting $\varepsilon = 0.1$ gives $t^2 \approx 1.2$, hence $t^* \approx 1.10$. (A more precise numerical evaluation gives $t^* \approx 0.74$ for 10\% relative error in the cost itself, not just the leading correction.)
\end{proof}

This crossover scale, being of order $\ln \varphi$, is a zero-parameter prediction: the strain at which classical elastic theory fails for quasicrystals is determined by the golden ratio, not by any material property.


\section{Open Problems and Falsification Criteria}
\label{sec:open}

\subsection{Falsification criteria}

Each sub-conjecture has a precise falsification criterion:

\begin{itemize}
\item \textbf{Sub-Conjecture A} is falsified by exhibiting a conservation-forced tiling whose eigenvalue is \emph{not} Pisot (e.g., a Salem number or transcendental).

\item \textbf{Sub-Conjecture B} is falsified by exhibiting a conservation-forced tiling that does \emph{not} admit a cut-and-project description, or whose height functions do not correspond to perpendicular-space coordinates.

\item \textbf{Sub-Conjecture C} is falsified by exhibiting a tiling where matching rules hold but cycle closure fails, or vice versa.

\item \textbf{Sub-Conjecture D} is falsified by exhibiting a Pisot substitution tiling that is \emph{not} conservation-forced (has $H^1_{PE} = 0$ or has no non-trivial PE 1-cocycles with non-PE primitives).

\item \textbf{Sub-Conjecture E} is falsified by exhibiting a conservation-forced tiling where the conservation rank $r \neq N - d$.
\end{itemize}

\subsection{Key open problems}

\begin{enumerate}
\item \textbf{Higher cohomology.} For Penrose tilings, $\Hcech^2(\Omega_P; \ZZ) \cong \ZZ^8$. Can the ledger framework produce natural 2-cochains generating $\Hcech^2$? A positive answer would extend the conjecture to ``conservation-forced at all degrees.''

\item \textbf{Non-substitutive tilings.} The conjecture is formulated for substitution tilings. Can it be extended to general cut-and-project tilings (which may lack a substitution rule)?

\item \textbf{The spectral gap conjecture.} Is $J(\lambda_\PF) \cdot (\lambda_\PF - |\lambda_2|) \geq C$ for some universal $C > 0$? Numerical evidence suggests this is \emph{false} as stated (the product approaches zero for eigenvalues near 1), but a modified version with an additional correction term may hold.

\item \textbf{Cost-weighted cohomology.} Define a ``$J$-weighted cohomology'' where cochains carry cost weights. Does this refinement distinguish tilings with the same ordinary cohomology but different cost structure?

\item \textbf{Three-dimensional verification.} Verify the conjecture for icosahedral tilings (3D quasicrystals with Pisot eigenvalue $\varphi$ and $N = 6$, $r = 3$).

\item \textbf{Physical observability.} The crossover scale~\eqref{eq:crossover} is a falsifiable prediction: in real quasicrystals, phason dynamics should deviate from quadratic elasticity at strain $\sim 1.5 \ln \varphi$. Can this be measured in diffuse scattering experiments?
\end{enumerate}


\section{Conclusion}

We have proposed a classification conjecture that unifies two structural correspondences between the cost-first ledger framework and aperiodic tiling theory. The central concept---conservation-forced aperiodic order---captures the idea that long-range quasicrystalline structure is a \emph{necessary consequence} of conservation with finite local resolution, not an accidental mathematical construction.

The conjecture, if true, would establish:
\begin{enumerate}
\item A \textbf{variational principle} for quasicrystal theory: the ``right'' tilings are cost-zero minimizers.
\item A \textbf{coherence hierarchy} among quasicrystals, with $\varphi$ at the unique minimum.
\item A \textbf{recognition-theoretic} foundation: the topology of tiling spaces encodes the gap between local observation and global structure.
\end{enumerate}

The conjecture is verified for all known 1D and 2D quasicrystalline substitution tilings. Five precise sub-conjectures with explicit falsification criteria provide a roadmap for future work.

\section*{Acknowledgments}
S.P.-G.\ thanks the developers of the cost-first ledger framework for ongoing collaboration. J.W.\ acknowledges support from the Recognition Science Research Institute.

\begin{thebibliography}{99}

\bibitem{PardoGuerraGolden2026}
S.~Pardo-Guerra,
Structural correspondences between cost-first ledgers and Penrose tilings,
submitted (2026).

\bibitem{PardoGuerraCohom2026}
S.~Pardo-Guerra,
Cohomological connections between the cost-first ledger framework and aperiodic tiling spaces,
submitted (2026).

\bibitem{PardoGuerraEtAl2026}
S.~Pardo-Guerra, M.~Simons, A.~Thapa, and J.~Washburn,
Coherent comparison as information cost: a cost-first ledger framework for discrete dynamics,
submitted (2026).

\bibitem{WashburnZlatanovic2026}
J.~Washburn and M.~Zlatanovi\'{c},
Uniqueness of the canonical reciprocal cost,
submitted (2026); arXiv:2602.05753v1.

\bibitem{Penrose1974}
R.~Penrose,
The role of aesthetics in pure and applied mathematical research,
\emph{Bull.\ Inst.\ Math.\ Appl.}\ \textbf{10} (1974), 266--271.

\bibitem{BaakeGrimm2013}
M.~Baake and U.~Grimm,
\emph{Aperiodic Order, Vol.~1: A Mathematical Invitation},
Cambridge University Press, 2013.

\bibitem{AndersonPutnam1998}
J.~E.~Anderson and I.~F.~Putnam,
Topological invariants for substitution tilings and their associated $C^*$-algebras,
\emph{Ergodic Theory Dynam.\ Systems}\ \textbf{18} (1998), 509--537.

\bibitem{KellendonkPutnam2006}
J.~Kellendonk and I.~F.~Putnam,
The Ruelle--Sullivan map for actions of $\mathbb{R}^n$,
\emph{Math.\ Ann.}\ \textbf{334} (2006), 693--711.

\bibitem{Salem1963}
R.~Salem,
\emph{Algebraic Numbers and Fourier Analysis},
Heath Mathematical Monographs, Boston, 1963.

\bibitem{GrunbaumShephard1987}
B.~Gr\"{u}nbaum and G.~C.~Shephard,
\emph{Tilings and Patterns}, W.~H.~Freeman, New York, 1987.

\bibitem{Sadun2008}
L.~Sadun,
\emph{Topology of Tiling Spaces},
University Lecture Series vol.~46, American Mathematical Society, 2008.

\bibitem{Forrest2002}
A.~H.~Forrest, J.~R.~Hunton, and J.~Kellendonk,
\emph{Topological Invariants for Projection Method Patterns},
Mem.\ Amer.\ Math.\ Soc.\ \textbf{159}, no.~758, 2002.

\bibitem{Shechtman1984}
D.~Shechtman, I.~Blech, D.~Gratias, and J.~W.~Cahn,
Metallic phase with long-range orientational order and no translational symmetry,
\emph{Phys.\ Rev.\ Lett.}\ \textbf{53} (1984), 1951--1953.

\bibitem{Levine1985}
D.~Levine, T.~C.~Lubensky, S.~Ostlund, S.~Ramaswamy, P.~J.~Steinhardt, and J.~Toner,
Elasticity and dislocations in pentagonal and icosahedral quasicrystals,
\emph{Phys.\ Rev.\ Lett.}\ \textbf{54} (1985), 1520--1523.

\end{thebibliography}

\end{document}
